Fig.~\ref{fig:fp-audio} and Table~\ref{tab:fp-audio} report the audio track. As in the image track, the corpus contained duplicates: FMA carries some recordings under two track identifiers, which all three recording-level systems matched perfectly. We removed the \val{fp:audio/dedup/neg} negative sources that were copies of indexed tracks and counted a duplicate of a registered track as the same asset (\val{fp:audio/dedup/items} registered tracks had one), using the same any-of-three rule as for images. Without this step every method's calibrated detection falls by half or more (e.g. Chromaprint \val{fp:audio/Chromaprint/nodedup/tpr} instead of \val{fp:audio/Chromaprint/TPR@query@0.01}); we report the corrected values. With \val{fp:count/audio/neg_cal} calibration queries the pair-level targets below $10^{-6}$ are not resolved, so we compare at the query-level 1\% threshold.

The neural music fingerprint NMFP detects the most transformed queries (\val{fp:audio/NMFP-triplet/TPR@query@0.01/pct}\%), followed by Chromaprint (\val{fp:audio/Chromaprint/TPR@query@0.01/pct}\%) and audfprint (\val{fp:audio/audfprint/TPR@query@0.01/pct}\%); the ISCC Audio-Code reaches \val{fp:audio/ISCC-Audio-256/TPR@query@0.01/pct}\% at 256 bits and \val{fp:audio/ISCC-Audio-64/TPR@query@0.01/pct}\% at its default 64 bits, and the semantic CLAP embedding \val{fp:audio/CLAP/TPR@query@0.01/pct}\%. The families fail on different transformations (Fig.~\ref{fig:fp-audio}a). Codecs, resampling and moderate noise are harmless to all of them. Excerpts separate the systems that align in time from those that summarise a whole file: Chromaprint, audfprint and NMFP find a 5\,s excerpt of a 30\,s track in \val{fp:audio/Chromaprint/att/excerpt5/pct}\%, \val{fp:audio/audfprint/att/excerpt5/pct}\% and \val{fp:audio/NMFP-triplet/att/excerpt5/pct}\% of cases, while the ISCC Audio-Code, a single SimHash over the whole Chromaprint vector, finds \val{fp:audio/ISCC-Audio-64/att/excerpt5/pct}\%. Loud babble from other music (0\,dB) defeats the whole-file codes and CLAP, and most Chromaprint matches (\val{fp:audio/Chromaprint/att/babble_0db/pct}\%), while the landmark system audfprint (\val{fp:audio/audfprint/att/babble_0db/pct}\%), whose sparse spectral peaks survive mixing, and NMFP (\val{fp:audio/NMFP-triplet/att/babble_0db/pct}\%) keep most of them. Babble results depend on which background tracks are mixed in: the two runs of this benchmark drew different backgrounds (Section~\ref{sec:fp-method}), and NMFP's rate at 0\,dB was about half as high with the first set. Pitch and speed changes of a few percent defeat every recording-level fingerprint (Chromaprint \val{fp:audio/Chromaprint/att/pitch+2/pct}\% at $+2$ semitones, NMFP \val{fp:audio/NMFP-triplet/att/pitch+2/pct}\%), and only CLAP, which describes content rather than the recording, survives them (\val{fp:audio/CLAP/att/pitch+2/pct}\%) at the cost of failing under noise (\val{fp:audio/CLAP/att/noise_10db/pct}\% at 10\,dB SNR). The simulated re-recording is handled well by NMFP (\val{fp:audio/NMFP-triplet/att/rerecord/pct}\%) and audfprint (\val{fp:audio/audfprint/att/rerecord/pct}\%) and poorly by the others.

Speech behaves like music for the recording-level systems (Recall@1 on 10\,s excerpts: Chromaprint \val{fp:audio/Chromaprint/speech/att/excerpt10/pct}\% on speech and \val{fp:audio/Chromaprint/music/att/excerpt10/pct}\% on music), and ISCC does better on speech than on music excerpts (\val{fp:audio/ISCC-Audio-256/speech/att/excerpt10/pct}\% against \val{fp:audio/ISCC-Audio-256/music/att/excerpt10/pct}\% at 256 bits). NMFP is the only reference-tier audio method; it is GPL-3.0 and was trained on FMA, whose training tracks we could not exclude from our FMA-large corpus by identifier, so its lead may be partly in-distribution.

\begin{figure*}[!t]
\centering
\includegraphics[width=\textwidth]{fp_fig6_audio}
\caption{(a) Recall@1 of each audio method under each transformation (\val{fp:corpus/audio/reg} registered items, \val{fp:count/audio/nref} in the index). (b) Detection against the pair-level false-positive rate; unresolved targets are omitted.}
\label{fig:fp-audio}
\end{figure*}

\begin{table*}[!t]
\centering
\scriptsize
\caption{Audio retrieval over all attacks, and Recall@1 under selected attacks. TPR$_{q,1\%}$: detection at the threshold set for a 1\% query-level false-match rate on the calibration negatives; FPR$_q$ (test): the rate that threshold actually gave on the held-out negatives. TPR$_{p,10^{-6}}$: the strictest pair-level target that 4053 calibration queries against 6350 references resolve; stricter targets are not reported. Intervals: 95\% bootstrap over sources.}
\label{tab:fp-audio}
\setlength{\tabcolsep}{3pt}
\begin{tabular}{lccccccccc}
\toprule
Method & R@1 & $\mu$AP & TPR$_{q,1\%}$ & FPR$_{q}$ (test) & TPR$_{p,10^{-6}}$ & Music 10\,s exc. & Speech 10\,s exc. & Pitch +2 & Re-record \\
\midrule
Chromaprint & 0.805 & 0.817 & 0.704 \tiny[0.68, 0.72] & 0.016 \tiny[0.00, 0.04] & 0.688 \tiny[0.66, 0.71] & 1.000 & 1.000 & 0.006 & 0.451 \\
audfprint & 0.851 & 0.851 & 0.660 \tiny[0.53, 0.74] & 0.008 \tiny[0.00, 0.03] & 0.599 \tiny[0.46, 0.72] & 0.982 & 1.000 & 0.071 & 0.852 \\
ISCC Audio-64 & 0.662 & 0.581 & 0.409 \tiny[0.21, 0.41] & 0.017 \tiny[0.00, 0.03] & 0.212 \tiny[0.00, 0.41] & 0.247 & 0.657 & 0.004 & 0.183 \\
ISCC Audio-256 & 0.749 & 0.658 & 0.454 \tiny[0.31, 0.51] & 0.019 \tiny[0.00, 0.04] & 0.321 \tiny[0.24, 0.47] & 0.555 & 0.882 & 0.006 & 0.342 \\
CLAP & 0.707 & 0.446 & 0.386 \tiny[0.33, 0.45] & 0.017 \tiny[0.01, 0.04] & 0.328 \tiny[0.27, 0.41] & 0.957 & 0.725 & 0.910 & 0.228 \\
NMFP\reftier & 0.905 & 0.937 & 0.862 \tiny[0.82, 0.88] & 0.018 \tiny[0.00, 0.04] & 0.846 \tiny[0.79, 0.88] & 1.000 & 1.000 & 0.057 & 1.000 \\
\bottomrule
\end{tabular}

\end{table*}

